\documentclass[12pt,a4paper]{article}

% ================================================================
% PAGE SETUP 
% ================================================================

\usepackage[a4paper,left=2cm,right=2cm,top=2cm,bottom=2cm]{geometry}
\usepackage{setspace}
\usepackage{changepage}
\usepackage{titlesec}

\setcounter{secnumdepth}{0}
\titleformat{\section}
  {\normalfont\Large\bfseries\centering}
  {}{0pt}{}

% ================================================================
% TYPESETTING AND LANGUAGE
% ================================================================

\usepackage[T1]{fontenc}
\usepackage{enumitem}
\usepackage{newtxtext,newtxmath}
\usepackage[ngerman]{babel}
\usepackage{csquotes}
\usepackage[normalem]{ulem}
\newenvironment{longquote}
{%
    \vspace{-1em}
    \begin{adjustwidth}{2cm}{2cm}
        \singlespacing
        \fontsize{10}{12}\selectfont
    }
{%
    \end{adjustwidth}
}

% ================================================================
% FOOTNOTES AND BIBLIOGRAPHY   
% ================================================================

\usepackage[
    backend=biber,
    style=deutschstellenbosch,
    autocite=footnote
]{biblatex}

\addbibresource{quellen.bib}

\makeatletter
\renewcommand{\footnotesize}{%
    \def\baselinestretch{1}%
    \@setfontsize\footnotesize{10}{12}%
}
\makeatother

\usepackage[hidelinks,hyperfootnotes=true]{hyperref}

% ================================================================
% METADATA
% ================================================================

% Replace these values with your own content.
\newcommand{\CoverCourse}{Deutsch 358}
\newcommand{\CoverInstructor}{Professor Ingo Müller}
\newcommand{\CoverDate}{den 30. Oktober 2026}
\newcommand{\CoverDepartment}{Department of Modern Foreign Languages}
\newcommand{\CoverStudent}{Barnard, Quin Charl \#27902277}
\newcommand{\CoverEmail}{27902277@sun.ac.za}

\begin{document}
\hypersetup{pageanchor=false}

% ================================================================
% COVER PAGE
% ================================================================

\pagenumbering{gobble}
\begin{titlepage}
\thispagestyle{empty}

\vspace*{9cm}

{\raggedleft
{\fontsize{30}{35}\selectfont \texttt{/titel/}\par}
\par}

\vspace{0.76cm}

{\raggedleft
    {\itshape\fontsize{18}{22}\selectfont \texttt{/untertitel/}\par}
\par}

\vspace{1.1cm}

\begin{raggedleft}
   {\fontsize{18}{21}\selectfont \CoverCourse\par}
    \vspace{0.07cm}
    {\fontsize{11}{14}\selectfont
      \CoverInstructor\par
      \CoverDate\par
      \CoverDepartment\par}
\end{raggedleft}

\vspace{1.4cm}

\begin{raggedleft}
    {\fontsize{15}{18}\selectfont \CoverStudent\par}
    \vspace{0.1cm}
    {\fontsize{10}{12}\selectfont \CoverEmail\par}
\end{raggedleft}
\end{titlepage}

% ================================================================
% TABLE OF CONTENTS
% ================================================================

\clearpage

\hypersetup{pageanchor=true}
\pagenumbering{arabic}
\setcounter{page}{1}
\tableofcontents

\clearpage

% ================================================================
% BODY
% ================================================================

\onehalfspacing
\setlength{\parskip}{0.5\baselineskip}

\section{Einleitung}
%
Im Jahr 1912 erscheint Franz Kafkas Band \textit{Betrachtung}\footnote{Kafkas Sammlung \textit{Betrachtung} wird im Folgenden zitiert nach: \cite{kafka}}, der Erzählungen enthält, die sowohl von Angst, Scham und Schuld als auch von weniger drastischen Affekten geprägt sind. Affekte wie Unbehagen, Unsicherheit\footnote{Es muss zwischen Unsicherheit als Zweifel an den Tatsachen und als Gefühl der Gefahr unterschieden werden. Der vorliegende Aufsatz bezeichnet erstere als „epistemologische Unsicherheit“ und letztere als „affektive Unsicherheit“.} und Verdruss sind, insofern sie beim Subjekt einen Wunsch nach eine Änderung der Lage auslösen, den mehr hervorstechenden Affekten ähnlich. Die drastischeren Affekte sind schon vielfach analysiert wurden, daher lohnt sich eine Analyse, die versucht, sich mit diesen weniger beachteten Affekten auseinanderzusetzen.\par
%
Die beiden längsten Erzählungen, nämlich \textit{Kinder auf der Landstraße} und \textit{Unglücklichsein} sind zu diesem Zweck besonders geeignet, da sie sehr detailliert darstellen, wie alltägliche Ereignisse von den Protagonisten erlebt werden und wie sie diese aus dem Gleichgewicht bringen. Es gibt eine weitere Ähnlichkeit zwischen den erwähnten Geschichten: es ist jeweils die Rede von einem Protagonisten, der auf jemanden wartet. Als die erwartete Person zum Vorschein kommt, streiten sich aber die Figuren, und sie sind sich nicht darüber einig, was zwischen ihnen geschieht und geschehen ist. Wie genau diese Verunsicherung sich gestaltet und auf den Leser wirkt, kann anhand der Affekte wie Unsicherheit, Verdruss und Unbehagen untersucht werden. Wegen den Ähnlichkeiten, lassen sich die Geschichten vergleichen, um die Frage der Verunsicherung erfolgreich erforschen zu können.\par 
%
\clearpage
%
\section{Kinder auf der Landstraße}
%
Der Satzteil \textit{„[…] daß es eine Schande war“} (6) als eine Gruppe Arbeiter am Anfang der Geschichte vorbeikommen, überrascht den Leser. Das Gelächter steht normalerweise für die Heiterkeit, weiterhin ist es üblich von irgendeinem Benehmen zu sagen „dass es eine \uline{Freude} war.“ Das Ereignis über dem sich der Protagonist hier fast beschwert ist alltäglich, doch findet er es so widrig, dass er einen festen Ausdruck umkehrt. Der Protagonist geht aber nicht weiter auf die Irritation ein: es wird hier hier nur gezeigt, dass sich der Protagonist wegen Alltagsituationen irritiert fühlt, was im Folgenden häufig thematisiert wird. \par
%
\begin{longquote}
    „Dann flogen Vögel wie sprühend auf, ich folgte ihnen mit den Blicken, sah, wie sie in einem Atemzug stiegen, bis ich nicht mehr glaubte, daß sie stiegen, sondern daß ich falle, und fest mich an den Seilen haltend aus Schwäche zu schaukeln anfing.“ (6)
\end{longquote}

Hier ist die Tätigkeit des Erzählers an sich zwar eine Alltägliche, doch bietet der Erzähler eine interessante Perspektive. Bewegung wird als unmittelbar relativ dargestellt, das heißt, der Bezugsrahmen kann notwendigerweise nicht absolut festgelegt werden. Der Perspektivwechsel ist für den Protagonisten herausfordernd, also muss er sich \textit{„aus Schwäche“} (6) an den Seilen halten, wobei er zu schaukeln beginnt. Indem er zunächst die Höhe der Vögel als zunehmend betrachtet, dann seine eigene Höhe als abnehmend, kehrt sich seine Perspektive völlig um. Dieser Unzuverlässigkeit seiner eigenen Perspektive schwächt den Protagonisten und er versucht sich statt auf sich selbst auf die Seilen zu verlassen. Diese sind aber gleich wie ihm mit der Schaukel verbunden und stehen daher nicht fest. Der eine schwankende Boden wird fast wortwörtlich gegen einen anderen getauscht, jedoch beginnt er bald stärker zu schaukeln. Ursprünglich war das Schaukeln lediglich die Folge seines Unbehagens, und jetzt tauchen zusammen sowohl ein angenehmer Wind als auch „zitternde Sterne [statt der fliegenden] Vögel“ auf. Offensichtlich geht es dem Erzähler besser als früher. Horstkotte führt dem Leser die „Spannung zwischen Mobilität und Immobilität“\autocite[128]{horstkotte2014} vor Augen, wo Immobilität mit sinnlicher Wahrnehmung assoziiert ist. Weil der Protagonist sich aktiv Mühe machen muss um zu schaukeln, und sich noch immer auf diesen instabilen Boden befindet, kann der Leser diese Lage im Text als eine Akzeptanz der Relativität interpretieren. Die Bewegung und die Aktivität kreiert diese Akzeptanz, wo Immobilität und Passivität für den Protagonisten kaum ertraglich sind. Laut Fuchs wird hier mit der „Inszenierung von Beschleunigung und Geschwindigkeitsrausch“\autocite[59]{fuchs} experimentiert. Di Mauro nennt der Rausch an dieser Stelle „ein kindliches, nicht kontrollierbares Nicht-Aufhören-Wollen“.\autocite[56]{mauro} Der Rausch und das dazugehörige Experimentieren entstehen nur insofern als sie die epistemologische Unsicherheit des Perspektivwechsels auflöst. Diese Auflösung ist warum das Ich, wie Di Mauro sagt, nicht aufhören will. Die bloße Wahrnehmung, die früher Angst erregte, wird durch den Rausch und das Unbehagen zum Mittel der positiven Änderung — obwohl die Wahrnehmung stets unzuverlässig sein kann, wie die Vögelszene zeigt.\par 
%
Als nächstes wird der Protagonist vom Fenster aus gefragt, der Leser bekommt aber die Frage, die den Protagonisten staunen lässt, nicht mit. Der Fragende wird mit einer Gebirge, die oft in der Literatur die Angst oder das Schreck darstellt, verglichen. Daher kann der Leser ableiten, dass die Frage oder sogar der Fragende selbst beim Protagonisten Furcht erregt. Der Fragende wird auch mit der \textit{„[bloßen] Luft“} (6) verglichen, als ob er endlos und daher schwierig zu verstehen ist. Bisher hat der Protagonist lediglich seine Gegenstände betrachtet und er fand schon das gewissermaßen unangenehm. Als ihn plötzlich jemand herbeizieht, weiß er überhaupt nicht was, er tun oder sagen soll. Hier wird „das Verhältnis von Wahrnehmung und Subjektivität“\autocite[103f]{fromm} dargestellt. Das Ich weiß nicht, wie es sich äußern soll, doch ist seine Antwort dem Anderen nicht wichtig. Also betrachtet der Subjekt den Anderen, aber vermag es hier nicht tätig zu werden, wie er es früher fertigbrachte. Indem der Protagonist staunt, verköpert er die „Unmöglichkeit von Präsenz“.\autocite[104]{fromm} Der Protagonist findet wegen des Staunen unmöglich, ein normales Gespräch zu führen und in der Gegenwart ganz anwesend zu sein. Das merkt der Leser daran, dass diese Szene besonders vage und auf fast verschwommene Art beschrieben wird.\par
%
In der nächsten Szene ist das Gespräch auch ganz außergewöhnlich. Das andere Kind äußert die Gedanken des Lesers, indem es fragt, warum das Ich beim Aufstehen seufzt. Statt aber auf eine Antwort zu warten, fragt es sofort ob etwas schlimmes passieren ist. Als es fragt, ob alles wirklich verloren sei, denkt das Ich anscheinand nur die ablehnende Antwort, aber sagt es nicht. Dass das Ich schweigt, könnte wieder als „Unmöglichkeit von Präsenz“ interpretiert werden. Als das andere Kind zum Vorschein kommt, saß das Ich alleine im Dunkeln von Mücken umringt. Freilich ist diese keine angenehme Situation, und es muss nicht vergessen werden, dass das Ich hier passiv oder immobil ist. Als er mit dem anderen Kind vor dem Haus erscheint, beklagen sich die gesammten Kinder, dass die beiden nicht schnell genug gekommen sind. Es scheint als ob der Protagonist in dieser Gruppe nicht ganz willkommen ist, weil ihre Mitglieder eher unsympathisch sind, trotzdem macht er mit. Es ist hier von einer sozialen Verunsicherung die Rede, die der unklare Dialog darstellt. Der Leser weiß nicht genau, wer was sagt, weil die Sätze nur durch Bindestriche geteilt sind. Es ist aber deutlich, dass sich die Kinder streiten. Hier ist gleichzeitig eine affektive Unsicherheit und eine epistemologische Unsicherheit:
\vspace{-1em}
\begin{description}[
        style=multiline,
        leftmargin=12em,
        labelwidth=9em,
        labelsep=0.5em,
    ] 
        \item [Affektive Unsicherheit:] Die anderen Kinder sind dem Ich gegenüber unfreundlich, also die früher erwähnte soziale Unsicherheit.
        \item [Epistemologische Unsicherheit:] Das Ich weiß weder, was passieren wird, noch was genau die anderen Kinder mit ihren Aussagen meinen. Da jederzeit etwas gemeines gesagt werden kann, verschlimmert diese Unsicherheit die affektive.
\end{description}
\vspace{-1em}
%
Trotzdessen verlässt das Ich sein Zuhause. Es darf nicht vergessen werden, dass mit den Kindern gehen Aktivität heißt, wobei bleiben eine Passivität ist. Wie früher erwähnt, bevorzugt das Ich Aktivität oder Mobilität gegenüber Passivität beziehungsweise Immobilität, woher das Ich sich mit der unsympathischer Gruppe zusammensetzt.\par
%
Es ist möglich, den nächsten Teil als einen Rausch zu interpretieren, trotzdem scheint es als ob das zum Verdruss des Erzählers passiert. 
%
\begin{longquote}
    „Bald rieben sich unsere Westenknöpfe aneinander wie Zähne, bald liefen wir in gleichbleibender Entfernung, Feuer im Mund, wie Tiere in den Tropen. Wie Kürassiere in alten Kriegen, stampfend und hoch in der Luft, trieben wir einander die kurze Gasse hinunter und mit diesem Anlauf in den Beinen die Landstraße weiter hinauf.“ (7)
\end{longquote}
%
Dass diese Stelle im Text die Westenknöpfe mit Zähnen vergleicht, ist interessant. Der Leser wird dabei am Geräusch knirschender Zähne erinnert, was oft in der Literatur viele Bedeutnisse haben kann. Es ist möglich, die zusammenreibende Westenknöpfe als Metapher für „sich auf die Zunge beißen“ zu interpretieren. Eine solche Interpretation wird von einem folgenden Satzteil unterstützt: \textit{„[…] Feuer im Mund, wie Tiere in den Tropen.“} Ein metaphorischer Feuer wird zurückgehalten; vielleicht macht sich der Erzähler Mühe, irgendeinen Ärger nicht zu äußern — hier wird einen Verdruss geschildert. Dessen Auflösung ist schon mehrmals in dieser Geschichte aufgetaucht. Die Kinder liefen schon als die Westenknöpfe erwähnt wurden, daher das Zusammenreiben, also ist es nicht hier der Fall, dass das Ich passiv war und jetzt aktiv. Aber der Grad der Aktivität nimmt zu; die Kinder jagen einander und laufen weiter weg. Die Nennung \textit{„der alten Kriegen“} (7) wirkt als ob es irgendeinen bösen Feind gibt, vor dem die Kinder fliehen und sich verstecken müssen — hier ist von Angst und Furcht die Rede. Diese Interpretation belegt Schwarz: \textit{„die Kinder [entfliehen] im Schutze der Dunkelheit in die Landschaft“}.\autocite[528]{schwarz} Sie beschreibt, wie die Kinder in Kafkas Erzählungen fliehen, weil es Ihnen an Freiheit mangelt. Diese Meinung ähnelt Fromms Auffassung: \textit{„Im Zentrum der} Betrachtung \textit{steht demnach die Ausbildung von Individualität, bzw. die Ausbildung der ‚wahren Gestalt‘.“}\autocite[102]{fromm} Die Nähe der Anderen war zunächst die Quelle von Verdruss, der bald eine Furcht oder Angst wurde, vor der die ganze Gruppe eigentlich zu fliehen versucht. Die Mobilität und Aktivität wie früher besprochen als Auflösung irgendeiner Irritation wird hierdurch weniger metaphorisch und eher ein wortwörtlicher Ausweg.\par
%
Anscheinend war das Gleichnis mit Kürassiere mehr als nur Stilmittel: die Kinder trennen sich bald und kämpfen gegen einander. Es ist weder deutlich, wer zuerst kämpferisch wird, noch wer was sagt. Der Erzähler nicht glauben, dass diejenigen, die eigentlich zusammen eine Gruppe bilden sollen, einander jetzt feige nennen und verletzen wollen. Sowohl die affektive als auch die epistemologische Unsicherheit gehen weiter, da es tatsächlich an einen Grund für die Aufteilung mangelt, trotzdem stoßen die Kinder einander \textit{„vor die Brust“} (7). Die affektive wird bald aufgelöst, da der Kampf schnell vorbei ist und eine fast liebliche Lage folgt. Obwohl der Protagonist \textit{„fallend und freiwillig“} (7) kapituliert, wundert es noch den Leser, warum sich die Kinder gestritten haben und vor allem warum der Protagonist so schnell aufgegeben hat. Gosetti-Forencei ist der Meinung, das Streben nach Wissen scheitere immer Kafkas Protagonisten.\autocite[1]{gosetti2021} Gerade an dieser Stelle akzeptiert der Protagonist die epistemologische Unsicherheit, weil es ihm bequemer ist, was der Leser wieder an der Bewegung merkt.\par
%
Später in der Erzählung taucht ein Gleichnis auf, das sehr schwierig ist zu deuten: \textit{„Wenn man seine Stimme unter andere mischt, ist man wie mit einem Angelhaken gefangen.“} (9). Der Protagonist wollte singen und zwar als Mitglied einer Gruppe. Fuchs beschreibt wie der Beschleunigungs- und Geschwindigkeitsrausch hier wieder zum Vorschein kommen,\autocite[see 67]{fuchs}, doch bedeutet vom Angelhaken gefangen zu sein den Tod eines Fisches. Das Gleichnis kann daher als Zeichen fürs Gefühl der Hilfslosigkeit und Verlust der Freiheit des Willen interpretiert werden. Weil der Protagonist aber mit der Gruppe hierher gelaufen ist, genügt es nicht mehr einfach fortzulaufen. Er muss sich von den Anderen verabschieden, was er auch bald tut. Er begrüßt nur diejenigen, die ihm am nächsten sind, wobei niemand ihn ruft. Seine Lage als Außenseiter, der nicht ganz willkommen ist, und daher sein gesamtes Unbehagen im Text, wird hierdurch belegt. Wenn er wirklich als gleichberechtigt betrachtet wäre, da hätte sich jemand gefragt, wohin er geht. Er geht aber nicht nach Hause, wie der Satz \textit{„Es war schon Zeit.“} (9) den Leser vermuten lässt: er läuft wieder in den Wald und hat eine Stadt, wo die Leute nicht \textit{„schlafen […] weil sie nicht müde werden.“} (10) Wenn das Schlafen als Zustand der Bewusstlosigkeit betrachtet wird, wird es deutlich, dass der Protagonist danach sehnt, von der epistemologischen Unsicherheit frei zu sein. Das ist auch warum er sich von der Gruppe, in der man nie wissen kann, was als nächstes passieren wird, trennen musste.
\clearpage
%
\section{Unglücklichsein}
%
Diese Erzählung fängt mit einer Bewegung an. Dem Leser ist es noch nicht möglich zu wissen was \textit{„unertraglich geworden war“} (32) doch bekommt er mit, wie der Erzähler geschwind in sein Zimmer hinein kommt. Die beiden ersten Tatsachen, die der Leser in dieser Geschichte miterlebt sind irgendeine Irritation und rasche Bewegung, die wahrscheinlich beabsichtigt war, diese Irritation zu entkommen. Wie bei \textit{Kinder auf der Landstraße} fällt es dem Ich schwierig, Alltagsituationen zu vertragen: bloß ein Licht in der Gasse erschreckt ihn, aber als er schreit, tut er das nicht vor Furcht. Als er sagt, er schreie nur um den Schrei zu hören, wundert sich der Leser, was er damit meint. Es ist möglich, dass das Ich schreit, weil es sich mit der eigenen Stimme beruhigen will. Auf den nächsten Teil des Satzes muss der Leser sehr genau achten: wenn ein Geräusch gehört wird, geschieht das nur, weil irgendeine Oberfläche es reflektiert. Das neue Ziel \textit{„im Grund des Spiegels“} (32) weist auf diese Tatsache hin. Die reflektierende Oberfläche nimmt aber auch notwendigerweise das Geräusch in gewissermaßen auf, und schwächt es dadurch. Der Protagonist sagt nichts shwäche der Schrei, was bedeuten muss, dass der Schrei überhaupt nicht reflektiert wird. Hierdurch wird das Scheitern der Selbstberuhigung bedingt, und damit auch das Fortdauern sowohl der Furcht als auch der Irritation. Di Mauro beschreibt, wie die Tür bei Kafka \textit{„eine Art unüberschreitbare Grenze, der gegenüber Kafka eine zögernde Haltung einnimmt“}\autocite[68]{mauro} darstellt. Diese Meinung belegt der nächsten Satzteil, als sogar die Pferde draußen physisch auf das Öffnen der Tür reagieren. Die offene Tür erweckt solche starke Gefühle, dass es sowohl dem Leser als auch den Pferden kaum ertraglich ist. Doch ist eine offene Tür gar nicht außergewöhnlich. Die offensichtliche Überreaktion geht weiter als der Text, gleich wie \textit{Kinder auf der Landstraße}, auf einen Krieg verweist, was eine Gefahr suggeriert. Der Protagonist fühlt sich irgendwie bedroht und es gelingt ihm nicht, sich selbst zu beruhigen, also empfindet er sehr stark sowohl eine affektive als auch eine epistemologische Verunsicherung. Die zwei Texte unterscheiden sich dadurch, dass der Protagonist von bei \textit{Unglücklichsein} nicht die Gelegenheit hat, diese Verunsicherung zu entkommen, sondern sie dauert einfach fort.\par
%
Der Text ähnelt bald wieder \textit{Kinder auf der Landstraße} als einen \textit{„unmerklich schaukelnden Fußbodenbalken“} (32) erwähnt wird. Nochmals taucht der schwankende Boden als Zeichen für die epistemologische Unsicherheit auf, hier aber steht nicht das Ich auf den schwankenden Boden, sondern das Kind. Die epistemologische Verunsicherung wird hier auch von anderen Tatsachen symbolisiert, sowie die Lampe nicht im Korridor brennt und die \textit{„Dämmerung des Zimmers“} (32). Sogar der Satzteil \textit{„beruhigte sch aber unversehens mit dem Blick zum Fenster, vor dessen Kreuz der hochgetriebene Dunst der Straßenbeleuchtung endlich unter dem Dunkel liegen blieb“} (32) zeugt, nach Fromm, von Unsicherheit. Er schreibt, Gesten haben keine feste Bedeutungen, müssen daher interpretiert werden und seien folglich unmittelbar \textit{„[semantisch offen]“}.\autocite[111]{fromm} 
%
\clearpage
%
\noindent Wenn der Text mit \textit{Kinder auf der Landstraße} verglichen wird, kommt einen subtilen Unterschied zum Vorschein: 
\vspace{-1em}
\begin{description}[
        style=multiline,
        leftmargin=15em,
        labelwidth=12em,
        labelsep=0.5em
    ]
            \item [\textit{Kinder auf der Landstraße} :] Die epistemologische Unsicherheit wird häufig von Gegenständen, mit denen der Protagonist direkt umgeht, verkörpert — das Ich sitzt auf der Schaukel.
            \item [\textit{Unglücklichsein} :] Die epistemologische Unsicherheit nimmt der Protagonist wahr, und sie wird häufig von wahrgenommenen Gegenständen verkörpert — das Kind auf dem schaukelnden Fußbodenbalken.
\end{description}
\vspace{-1em}
%
Hierdurch wird wieder Fromms \textit{„Verhältnis von Wahrnehmung und Subjektivität“} dargestellt. Seiner Meinung nach, stellt \textit{Betrachtung} \textit{„die Frage nach den Möglichkeiten der Selbstaussage in der Moderne am Beispiel der Wahrnehmung“}.\autocite[104]{fromm} Das Ich wird in diesem Text, wegen der erwähnten Verunsicherung, fast völlig stumm, weil die Subjektivität seiner Wahrnehmung ihn so stark überrascht. Er hatte die falsche Erwartung, die eigene Stimme beim Schreien zu hören, und er vermag es nicht, das Benehmen des Kindes zu verstehen. Er sagt endlich \textit{„Guten Tag“} (32) als er seinen Rock nimmt. Dass er \textit{„so halb nackt“} (32) da stand, kann mit der epistemologischen Unsicherheit verknüpft werden: Adam und Eva verkleideten sich, als sie das verbotene Wissen beziehungsweise die Scham gelernt haben. Das heißt, sie waren nackt und so tief in der epistemologischen Unsicherheit, dass sie es nicht wussten. Der Protagonist erkennt hier so zu sagen sein Unwissen und löst die Verlegenheit auf indem er sich verkleidet und anfängt zu sprechen. Dadurch belegt er auch, dass er zur Selbstaussage fähig ist. Der Protagonist beruhigt und bestätigt sich selbst also hier, als es eine andere Person gibt, was er nicht früher mit dem Schreien konnte. Ein ganz alltägliches Ereignis hat die Verunsicherung verursacht, und sie wird wieder von einer anderen alltäglichen Handlung — der Gruß — aufgelöst. Die These, dass Äußerung das Ich der unangenehmen Affekten enthebt, bestätigt der Protagonist selbst: \textit{„Ein Weilchen lang hielt ich den Mund offen, damit mich die Aufregung durch den Mund verlasse.“} (32) Diese ist zwar eher im metaphorischen Sinn eine Äußerung, da er nichts sagt, doch öffnet er den Mund und wird somit von Gefühlen befreit, als hätte er gesprochen.\par
%
Als nächstes sagt der Protagonist, er habe \textit{„schlechten Speichel“} (32) in sich, und ihm zittern die Augenwimpern. Der Leser wurde hierher ableiten, dass das Ich sich unwohl fühlt, jedoch sagt der Protagonist, dass ihm \textit{„nichts las gerade dieser allerdings erwartete Besuch“} fehlt. Der Leser ist darum erstaunt, als er das Kind fragt, ob es sich sicher es, dass es hier sein will. Wenn der Protagonist den Besuch erwartet hat, und sich darauf gefreut hat, da würde man erwarten, dass er das Kind fröhlich willkommen heißt. Hier ist aber die Symbolik der Tür wie sie Di Mauro erläutert wichtig: sie erwähnt Freuds \textit{„Unheimlichen im Heim“}\autocite[71]{mauro} und wie es \textit{„die Wiederholung des Gleichen“} inszeniert. Die Tür ist schon als \textit{„unüberschreitbare Grenze“} bezeichnet, und es wird im Folgenden auch besprochen, in welchermaßen es von einer Wiederholung die Rede ist. Hauptsache ist, dass sich der Protagonist vor der offenen Tür fürchtet. Er versucht deswegen zuerst das Kind wegzumachen, durch die Frage ob es genau hier sein will, dann lädt er das Kind ein und fokussiert sich auf nichts als die Tür. Es scheitert dem Protagonisten, sich sozial normal zu benehmen, weil die Tür bei ihm Unbehagen auflöst. Das Kind empfindet offensichtlich dieses Unbehagen nicht, weil es sagt: \textit{„Ja, was ist denn? Was haben Sie? Meinetwegen kann das ganze Haus hereinkommen.“} Der Protagonist sagt später in der Erzählung, dass die Natur des Kindes die seine ist. Wie das möglich ist, wird anhand der Wiederholung besprochen.\par
%
Sehr viel später in der Erzählung, fast am Ende, hat der Protagonist sein Zimmer verlassen. Er beabsichtigt, spazieren zu gehen, nachdem er mit dem Nachbarn über dem Kind, das ein Gespenst sein soll, gesprochen hat. Er fühlt sich aber selbst verlassen und entscheidet sich dazu, schlafen zu gehen. Hier ist also von einem unertraglichen Affekt die Rede, der ihn zum Zimmer zwingt. Jetzt wird es deutlich, dass das Ende der Erzählung zusammen mit dem Anfang einen Kreis bildet. Das Gefühl der Verlassenheit, ist es, was unertraglich geworden ist, und die Tür ist eine \textit{„unüberschreitbare Grenze“} nicht weil er sie gar nicht überschreiten kann, sondern weil er sie immerhin überschreiten werden muss, und daher nie ein anderes Ziel als die Tür selbst erreichen kann. Die Tatsache, das die beiden Figuren die gleiche Natur haben, zeugt davon, dass sie die gleiche Person ist. Das Kind ist das Gespenst der Jugend von dem Protagonist. Das ist auch warum der Protagonist das Kind erwartet: er sehnt danach, wieder jung zu sein, aber er fühlt sich auch vor seiner Jugend unwohl. Er beschwert sich darüber, dass das Kind \textit{„wie verrückt [die Finger] an [seiner] Wand [reibt]“} (34f), doch gehört die Wand sowohl ihm als auch das Kind. Er kann es nicht akzeptieren, dass er gerade seine Kindheit konfrontieren muss und so durch die Gelegenheit hat, wieder die Jugend zu erleben, deswegen bleibt er als Subjekt in der Wahrnehmung stecken. Das Unbehagen und die gesamte Verunsicherung stärken die drastischeren Affekte, wodurch er in diesem Kreislauf gefesselt wird. 
\clearpage
%
\section{Schluss}
%
Es ist gedeutet worden, wie Affekte wie Verdruss, sowohl affektive als auch epistemologische Verunsicherung, Verlegenheit und andere, die mehr hervorstechenden Affekten unterstützen, um Krisen des Ich darzustellen. Der große Unterschied zwischen den beiden Texten, \textit{Kinder auf der Landstraße} und \textit{Unglücklichsein} steckt darin, wie sich die Protagonisten Mühe machen, diese Affekte und somit die Krisen zu entkommen. Die Bewegung und die Änderung sind stets wichtige Mittel zum Interpretation bei diesen Erzählungen, weil sie die Relativität und Subjektivität der Protagonisten betonen. Dass die Protagonisten tatsächlich die anderen Figuren erwartet haben, sich trotzdem unwohl fühlen, zeugt davon, dass es schwierig oder unmöglich ist, sich auf die untersuchten Affekte vorzubereiten.
% ================================================================
% BIBLIOGRAPHY
% ================================================================
\clearpage
\printbibliography[
    title={Literaturverzeichnis}
]

\end{document}
